\documentclass[11pt]{article}
\usepackage[margin=27mm]{geometry}
\usepackage[T1]{fontenc}
\usepackage{lmodern,amsmath,amssymb,amsthm,mathtools,microtype,hyperref}
\hypersetup{hidelinks}
\newcommand{\ddc}{dd^c}
\newcommand{\MA}{\operatorname{MA}}
\newcommand{\PSH}{\operatorname{PSH}}
\newtheorem{theorem}{Theorem}
\newtheorem{lemma}[theorem]{Lemma}
\title{Fixed spherical measures and exact Monge--Amp\`ere\newline
equations on Hermitian manifolds}
\author{}\date{9 October 2026}
\begin{document}\maketitle
\noindent\textbf{Scope.} The spherical sufficient-condition theorem and
the regularity lemma below have undergone independent reconstruction and
final mathematical review. Necessary classifications, optimal coefficients
and full original S15 completion are not asserted.
\begin{abstract}
We give a geometric sufficient test on singular measures for scaled
subsolutions, normalized exact homogeneous equations and exact exponential
equations on every compact Hermitian manifold. Fix a holomorphic ball and
the probability obtained by placing geometric weights on its concentric
shrinking spheres. Every finite nonzero measure with bounded density
against this predetermined probability is dominated by the product of an
explicit H\"older continuous potential. We use the analytic estimates of
Ko\l odziej--Nguyen, with a continuous-density diagonal and a coordinated
regularization parameter, for the homogeneous equation. The bounded-reference
theorem of Alehyane--Lu--Salouf gives the exponential equation. No positive
lower mass or finite upper mass of the background is required.
\end{abstract}

\section{The individual geometric test}
Let $(X,h)$ be any connected compact Hermitian manifold of complex
dimension $n\geq2$. Use $\ddc=i\partial\bar\partial/(2\pi)$ and
$d^c=i(\bar\partial-\partial)/(4\pi)$, genuine canonical
quasi-psh representatives, and the bounded smooth-background corrected
Bedford--Taylor product. Fix a holomorphic coordinate ball about $p$,
write $R=|z|$, and choose $0\leq\chi\leq1$ supported there and equal
to one near a smaller closed ball. Fix $r_j=r_0 2^{-j}$, $j\geq1$,
so small that every sphere $\{R=r_j\}$ lies in that smaller ball.
Write $\sigma_j$ for its normalized Euclidean surface probability,
transported to $X$ and extended by zero. Define, \emph{before the input
measure or exponent},
\[
 S_n=\sum_{j\geq1}r_j^{n/2},\qquad
 \rho=\frac1{S_n}\sum_{j\geq1}r_j^{n/2}\sigma_j.
\]
For a finite nonzero positive measure $\nu$, the test is
\[
 M_\rho(\nu)=\left\|\frac{d\nu}{d\rho}\right\|_{L^\infty(\rho)}<\infty,
\]
with value infinity if absolute continuity fails. It is a test on the
given measure against an explicit geometric measure, with no quantifier
over solutions. Changing coordinates while transporting the same fixed
measure changes neither the inequality nor the essential supremum.
Alternatively, the existence of \emph{some} holomorphic ball and radius
with this stated spherical test defines a coordinate-independent
geometric sufficient condition on $\nu$; no potential appears in it.

\begin{theorem}\label{thm:main}
The fixed $\rho$ is a singular non-pluripolar probability. There are
an explicit $k>0$ and a genuine $h$-psh $u_\rho\in C^{0,1/2}(X)$,
normalized by $\sup_Xu_\rho=0$, such that
\[
 \MA_h(u_\rho)\geq k\rho,\qquad
 0<\MA_h(u_\rho)(X)<\infty.
\]
For every finite nonzero $\nu$ with $M_\rho(\nu)<\infty$, the same
reference satisfies $\MA_h(u_\rho)\geq k\nu/M_\rho(\nu)$.
There exist a genuine H\"older continuous $h$-psh $w$ and a finite
constant $c>0$ such that
\[
 \sup_X w=0,\qquad\MA_h(w)=c\nu,
 \qquad0<\MA_h(w)(X)<\infty.
\]
For every specified $\lambda>0$ there is a unique bounded genuine
$h$-psh $U_\lambda$ such that
\[
 \MA_h(U_\lambda)=e^{\lambda U_\lambda}\nu,
 \qquad0<\MA_h(U_\lambda)(X)<\infty.
\]
Uniqueness in the last assertion is within bounded potentials. Its fixed
scalar allows no further supremum-zero normalization. The reference
measure and the test are independent of $\lambda$.
\end{theorem}

\section{The explicit H\"older reference}
Set
\[
 \psi_j=\chi\log\max(R^2,r_j^2)+(1-\chi),
\]
equal to one outside the coordinate ball. These functions are continuous
and bounded. Where $\chi=1$ they are psh. On the cutoff transition
$R$ is bounded away from zero and exceeds every $r_j$, so every
$\psi_j$ equals the \emph{same} fixed smooth function there.
Outside the cutoff they are constant. Thus one fixed finite $K\geq1$
satisfies $\ddc\psi_j\geq-Kh$ for all $j$. This is a global
uniform Hessian estimate, independent of the shrinking radius.
Put
\[
 Z=2K\sum_j\sqrt{r_j},\qquad d_j=\frac{\sqrt{r_j}}Z,
 \qquad u_0=\sum_jd_j\psi_j.
\]
Since $\sum d_j=1/(2K)$ and
$\sum d_j\|\psi_j\|_\infty<\infty$, the series converges uniformly,
and its genuine limit satisfies $h+\ddc u_0\geq h/2$.

Here is the actual H\"older estimate. In the inner ball, for two points
of Euclidean distance $t$, the reverse triangle inequality gives
\[
 |\log\max(|z|^2,r_j^2)-\log\max(|z'|^2,r_j^2)|
 \leq2\log(1+t/r_j).
\]
For $r_j\geq t$, the weighted sum is bounded by
$C t\sum_{r_j\geq t}r_j^{-1/2}\leq C\sqrt t$.
For $r_j<t$, write the radii as a geometric tail below $t$; its sum
is bounded by
$C\sqrt t\sum_{l\geq0}2^{-l/2}(1+l)\leq C\sqrt t$.
On the transition and exterior, the radius-dependent logarithm has
disappeared and the sum is a fixed smooth function. A finite chart
cover therefore proves $u_0\in C^{0,1/2}(X)$.

In $\mathbb C^n$ the actual bounded product is
\begin{equation}\label{eq:sphere}
 (\ddc\log\max(R^2,r^2))^n=\sigma_r.
\end{equation}
Indeed the potential is constant inside, while its exterior Hessian
has rank $n-1$. Rotation invariance leaves only a constant surface
measure at the sphere. Its total mass is
$\int_{R=r}d^c\log R^2\wedge(\ddc\log R^2)^{n-1}=1$:
the Hopf circle integral is one and the induced Fubini--Study
$(n-1)$-volume on $\mathbb P^{n-1}$ is one in this normalization.
The identity can equivalently be obtained by smooth convex radial
regularization and bounded-product continuity \cite{BT}.
This bounded local psh product gives zero mass to pluripolar sets.
Consequently each $\sigma_j$, and then their countable positive sum
$\rho$, is non-pluripolar. The countable union of spheres has zero
smooth volume, so $\rho$ is singular. This is countable additivity,
not preservation of non-pluripolarity under an arbitrary weak limit.

On the smaller ball all $\ddc\psi_j$ are closed positive currents.
For a finite sum the corrected bounded product of
$h+\sum_{j\leq N}d_j\ddc\psi_j$ expands with smooth $h>0$;
all terms are positive. Retain any fixed finite set of diagonal terms
$d_j^n\sigma_j$. Uniform convergence of the local corrected potentials
gives convergence of the actual products \cite{BT}. Pass against
positive tests and exhaust that finite set. Thus globally
\[
 \MA_h(u_0)\geq\sum_jd_j^n\sigma_j=k\rho,
 \qquad k=\frac{S_n}{Z^n}>0.
\]
This only uses local closed positive currents plus a smooth positive
background; it imports no arbitrary global nonclosed mixed-product
theorem. Boundedness and the local bounded-product estimates on a finite
cover make the total product finite, even when $v_+(h)=\infty$.
Its domination gives strict positivity. Subtracting $\sup_Xu_0$
defines $u_\rho$ without changing its product.

\section{Two exact range applications and their scope}
The precise range hypothesis in Ko\l odziej--Nguyen
\cite[Definition 1.1 and Theorem 1.3]{KN} is domination of a finite
nonzero measure by the product of a H\"older continuous global
$h$-psh reference. Our construction supplies it. We record the required
approximation and regularity details below, using their analytic estimates.
This also avoids applying a fixed Borel density to weakly converging
smooth reference products.

Write $\alpha=\MA_h(u_\rho)$ and $\nu=C_0g\alpha$, with $0\leq g\leq1$.
Choose continuous $g_l\in[0,1]$ tending to $g$ in $L^1(\alpha)$.
Then $g_l^{1/n}\to g^{1/n}$ in $L^1(\alpha)$. By
\cite[Lemma 2.3]{KN}, smooth genuine reference approximants $v_j$ tend
uniformly to $u_\rho$ with a common smaller H\"older norm. Choose a
diagonal $j(l)$ preserving this convergence and the integrals of
$g_l^{1/n}$ and of $g_l$ against the first $l$ elements of a dense
continuous test family, and set $\widehat v_l=v_{j(l)}$. Thus
\[
 \nu_l=C_0g_l\MA_h(\widehat v_l)\rightharpoonup\nu,\qquad
 \inf_{l\gg1}\int_Xg_l^{1/n}\MA_h(\widehat v_l)>0.
\]
These data have continuous smooth-volume densities and positive mass.
The continuous $L^p$ theorem \cite[Theorem 0.1]{KN13} gives normalized
continuous $h$-psh $\phi_l$ with $\MA_h(\phi_l)=c_l\nu_l$, $c_l>0$.
The mixed inequality and weighted CLN estimate
\cite[Proposition 1.1 and Lemma 1.9]{Ngu} give
\[
 c_l^{1/n}C_0^{1/n}\int_Xg_l^{1/n}\MA_h(\widehat v_l)
 \leq\int_Xh_{\phi_l}\wedge h_{\widehat v_l}^{n-1}\leq C,
 \qquad h_a=h+\ddc a.
\]
Hence $c_l$ is bounded above. The common H\"older-reference capacity
bound \cite[Lemma 2.1 and Corollary 2.2]{KN} consequently applies to every
$\MA_h(\phi_l)$. Modified comparison with the zero competitor and
\cite[Theorem 5.3 and Lemma 5.4]{KN13} give, for one fixed small $s>0$,
\[
 \operatorname{cap}_h\{\phi_l<\inf_X\phi_l+s\}
 \leq Cc_l\nu_l(X)/s^n,\qquad
 s\leq\kappa\bigl(\operatorname{cap}_h\{\phi_l<\inf_X\phi_l+s\}\bigr),
\]
where $\kappa(t)\to0$ as $t\to0$. This bounds $c_l$ away from zero.
The same capacity proof and the normalized sublevel-capacity estimate
used in \cite[Corollary 5.6]{KN13} give a uniform bound for $\phi_l$.
No uniform $L^p$ bound for the approximating densities is needed.

Extract $c_l\to c>0$ and $\phi_l\to\phi$ in $L^1(X)$.
We explain the two-index convergence needed for uniform stability.
For uniformly bounded genuine $h$-psh $a,b,r,t$, integration by parts,
local CLN and gradient energy bounds give
\begin{equation}\label{eq:transfer}
 \left|\int_X|a-b|\bigl(\MA_h(r)-\MA_h(t)\bigr)\right|
 \leq C\|r-t\|_\infty.
\end{equation}
Indeed use $|a-b|=2\max(a,b)-a-b$ and
$h_r^n-h_t^n=\ddc(r-t)\wedge\sum_{k=0}^{n-1}h_r^k\wedge h_t^{n-1-k}$.
In $\ddc(fT)$ retain all four terms
$\ddc f\wedge T+df\wedge d^cT-d^cf\wedge dT+f\ddc T$.
Their masses are uniformly bounded by CLN and Cauchy--Schwarz, since
$dh_r=dh_t=dh$ and the torsion is fixed and smooth
\cite[Propositions 1.4 and 1.5]{Ngu}. Bounded local regularization
justifies these operations; terms with vanishing coefficients are
omitted in low dimensions.
Each displayed mixed identity is a finite local expansion in closed
corrected currents and smooth forms, glued by their actual definitions.
For each fixed reference product, extract an $L^1$-convergent subsequence
using \cite[Lemma 2.1 and Corollary 2.2]{KN22}; a diagonal handles
the countably many fixed references. Uniform reference convergence and
\eqref{eq:transfer} then imply
\[
 \int_X|\phi_j-\phi_k|\MA_h(\widehat v_j)\longrightarrow0
 \quad\text{as }j,k\longrightarrow\infty,
\]
Domination transfers this to the solution products. Apply
\cite[Proposition 2.4]{KN} in both directions to obtain a uniform
Cauchy subsequence. Its continuous genuine limit satisfies
$\MA_h(\phi)=c\nu$ by bounded corrected-product continuity.
Lemma~\ref{lem:repair} below gives H\"older regularity. Its normalization
is already $\sup_X\phi=0$, and its mass is $c\nu(X)$.
We assert neither homogeneous uniqueness nor an optimal coefficient.

For the exponential equation use Alehyane--Lu--Salouf
\cite[Theorem 5.1]{ALS}. Its background is smooth semipositive and big;
$h>0$ is such a background. The domination gives the exact datum
\[
 \nu=f\MA_h(u_\rho),\qquad0\leq f\leq M_\rho(\nu)/k.
\]
The reference is bounded and the datum is finite and nonzero, as
required. That theorem supplies the unique bounded solution for each
positive exponent. In its bounded stability argument the rooftop
infimum is bounded and comparable to the zero model, the alternative
in Lemma 4.3; no positive-lower-mass assumption enters this application.
The separate relative-range Theorem 5.8, which has that additional
hypothesis, is not used.

Both applications preserve actual currents under a numerical operator
change: if $D_{\rm source}=s\ddc$, use the source reference
$u_\rho/s$. Transfer its output by multiplying by $s$; in the
exponential application use source exponent $s\lambda$.
The actual product and target are unchanged. The range theorems and
their analytic foundations remain explicitly imported results, not
consequences of the retained inequality alone.

For a background with $v_-(h)=0$ the same construction and both exact
equations still apply. It therefore gives a substantive singular datum
in that regime. No bounded-density truncation relative to smooth volume
is involved. If $v_+(h)<\infty$, each bounded product here is also
bounded by that upper mass, but the theorem does not require it.

The test is sufficient; failure of it makes no impossibility assertion.
It distinguishes the scaled inequality, the normalized homogeneous
equation and the fixed-scalar exponential equation by proving each with
its own normalization and range application. Necessary criteria,
unbounded uniqueness, prescribed masses, optimal coefficients and the
K3 bounded infimum are not conclusions. The abstract range theorems
are substantial prior work; no priority is claimed for them or for a
general classification. Original S15 completion must be assessed
against its unchanged criterion, not inferred from this author draft.

\section{A coordinated regularization parameter}
The following lemma supplies the regularity step using only
\cite[Proposition 2.4 and Lemmas 2.7--2.8]{KN}. It does not use the
final H\"older range theorem as a premise.
\begin{lemma}\label{lem:repair}
Let $\phi$ be a continuous genuine $h$-psh function with $\sup_X\phi=0$,
$\MA_h(\phi)=\mu$, and $0<\mu(X)<\infty$.
If $\mu$ is dominated by the product of a H\"older continuous global
reference, then $\phi$ is H\"older continuous.
\end{lemma}
\begin{proof}
Write $\mathcal R_t\phi$ for the normalized Chern regularization.
The imported stability and regularization estimates give exponents
$\gamma,\eta\in(0,1)$ such that, for continuous genuine $h$-psh $\psi\leq0$,
\[
 \sup_X(\psi-\phi)\leq C\|(\psi-\phi)_+\|_{L^1(\mu)}^\gamma,
 \qquad\|\mathcal R_\delta\phi-\phi\|_{L^1(\mu)}\leq C\delta^\eta.
\]
Moreover $\mathcal R_t\phi+Kt^2$ increases with $t$ and tends to $\phi$
at zero. For $A\geq1$ and $K>0$ fixed, the transform
\[
 F_{\delta,b}=\inf_{0<t\leq\delta}
 \left(\mathcal R_t\phi+Kt^2+Kt-b\log(t/\delta)\right)
\]
satisfies $h+\ddc F_{\delta,b}\geq-(Ab+2K\delta)h$.
Choose
\[
 q=\eta\gamma/2,\qquad\varepsilon=\delta^q,
 \qquad b=(\varepsilon-2K\delta)/A.
\]
For small $\delta$, $\varepsilon\leq1/2$ and
$\varepsilon/(2A)\leq b\leq\varepsilon/A$.
Set $T=(1-\varepsilon)F_{\delta,b}$. The curvature estimate gives
\[
 h+\ddc T=\varepsilon h+(1-\varepsilon)(h+\ddc F_{\delta,b})
 \geq\varepsilon^2h.
\]
Monotonicity and evaluation at $t=\delta$ imply
$\phi\leq F_{\delta,b}\leq\mathcal R_\delta\phi+2K\delta$.
The transform is continuous: the positive logarithmic penalty and
boundedness exclude $t$ near zero uniformly for each fixed $\delta,b$.
Let $M=\|\phi\|_\infty$ and choose $B\geq M+2K+1$.
Since $\mathcal R_\delta\phi\leq0$, $\psi=T-B\varepsilon$ is a continuous
genuine $h$-psh competitor with $\psi\leq0$. The crucial cancellation is
\[
 T-B\varepsilon-\phi
 \leq(1-\varepsilon)(\mathcal R_\delta\phi-\phi+2K\delta)
       +\varepsilon(M-B)
 \leq(1-\varepsilon)(\mathcal R_\delta\phi-\phi+2K\delta).
\]
Thus $\|(\psi-\phi)_+\|_{L^1(\mu)}\leq C\delta^\eta$ and stability yields
\[
 \sup_X(T-\phi)\leq C\delta^{\eta\gamma}+B\varepsilon\leq C_1\varepsilon.
\]
As $\phi\leq0$,
$0\leq F_{\delta,b}-\phi=(T-\phi+\varepsilon\phi)/(1-\varepsilon)
\leq2C_1\varepsilon$.
At an actual minimizer $t_0\in(0,\delta]$, monotonicity gives
\[
 b\log(\delta/t_0)\leq F_{\delta,b}-\phi\leq2C_1\varepsilon,
 \qquad t_0\geq\kappa\delta,\quad\kappa=e^{-4AC_1}>0.
\]
This $\kappa$ is uniform. Another application of monotonicity gives
$\mathcal R_{\kappa\delta}\phi-\phi\leq2C_1\delta^q$.
Rescale $r=\kappa\delta$ and use
$\mathcal R_r\phi-\phi\geq-Kr^2$ to conclude
\begin{equation}\label{eq:uniform-smoothing}
 \|\mathcal R_r\phi-\phi\|_\infty\leq C_2r^q.
\end{equation}

We give the final smoothing implication directly. In a relatively
compact coordinate chart the smooth exponential map has fiber derivative
the identity at zero \cite[Proposition 2.9]{Dem}. Its parameterized
local inverse rewrites the
kernel in \cite[(2.7)--(2.8)]{KN} as
\[
 \mathcal R_r\phi(z)=\int\phi(y)K_r(z,y)\,dy,\qquad
 K_r(z,y)=r^{-2n}a\bigl(z,(y-z)/r,r\bigr).
\]
Here $a$ has uniformly bounded first derivatives and compact support in
a fixed bounded set of its second variable. These properties follow
from the inverse function theorem, smooth metric coefficients and
the smooth cutoff; they hold on a finite chart cover. Differentiating
the kernel and integrating over its support gives
\[
 \int|\nabla_zK_r(z,y)|\,dy\leq C/r,
 \qquad\|\nabla\mathcal R_r\phi\|_\infty\leq CM/r.
\]
For points at small Hermitian distance $s$, integrate along a short
minimizing path and use \eqref{eq:uniform-smoothing}:
\[
 |\phi(x)-\phi(y)|\leq2C_2r^q+CMs/r.
\]
Choose $r=s^{1/(1+q)}$ in the valid smoothing interval. Boundedness
handles larger distances. Therefore $\phi\in C^{0,q/(1+q)}(X)$.
No optimal exponent is asserted.
\end{proof}

\begin{thebibliography}{9}
\bibitem{BT} E. Bedford and B. A. Taylor,
\emph{A new capacity for plurisubharmonic functions}, Acta Math.
149 (1982), 1--40.
\bibitem{KN} S. Ko\l odziej and N. C. Nguyen,
\emph{H\"older continuous solutions of the Monge--Amp\`ere equation
on compact Hermitian manifolds}, Ann. Inst. Fourier 68 (2018),
2951--2964, Definition 1.1 and Theorem 1.3,
\href{https://aif.centre-mersenne.org/item/10.5802/aif.3232.pdf}
{doi:10.5802/aif.3232}.
\bibitem{KN13} S. Ko\l odziej and N. C. Nguyen,
\emph{Weak solutions to the complex Monge--Amp\`ere equation on
Hermitian manifolds}, \href{https://arxiv.org/abs/1312.5491v2}
{arXiv:1312.5491v2}.
\bibitem{Ngu} N. C. Nguyen,
\emph{The complex Monge--Amp\`ere type equation on compact Hermitian
manifolds and applications}, \href{https://arxiv.org/abs/1501.00891v1}
{arXiv:1501.00891v1}.
\bibitem{KN22} S. Ko\l odziej and N. C. Nguyen,
\emph{Weak convergence of Monge--Amp\`ere measures on compact Hermitian
manifolds}, \href{https://arxiv.org/abs/2212.11550v1}
{arXiv:2212.11550v1}.
\bibitem{ALS} O. Alehyane, C. H. Lu and M. Salouf,
\emph{Monge--Amp\`ere equations with prescribed singularities on compact
Hermitian manifolds}, \href{https://arxiv.org/abs/2511.02339v1}
{arXiv:2511.02339v1}, Theorem 5.1 and Lemma 4.3.
\bibitem{Dem} J.-P. Demailly,
\emph{Regularization of closed positive currents of type (1,1) by the
flow of a Chern connection}, Proposition 2.9,
\href{https://www-fourier.univ-grenoble-alpes.fr/~demailly/manuscripts/regularization_flow.pdf}
{author-hosted manuscript}.
\end{thebibliography}
\end{document}
